% !Mode:: "TeX:UTF-8"
% !TEX program  = xelatex
\section{结论与展望}
本文围绕多 LoRA 适配器推理中的显存受限与频繁切换问题，提出了基于预测模型的适配器调度方法。该方法以 EMA 建模访问趋势，并结合队列压力、活跃状态与显存感知归一化构建统一效用函数，在在线约束下以低开销近似求解驻留集合优化问题。

实验结果表明，相比 FIFO、LRU 与 LFU 等被动策略，本文方法在命中率、延迟与淘汰次数等关键指标上具有一致优势，尤其在显存预算紧张和热点动态迁移场景中表现更优。

未来工作可从三个方向继续推进：一是引入更强的学习式预测器，提升突发负载下的预测精度；二是将调度目标与端到端服务质量（如 tail latency SLO）联合建模；三是扩展至多 GPU 与分布式部署环境，研究跨设备的适配器协同驻留与迁移机制。
